\section{讨论}

\subsection{PC-RCO的鲁棒性来源}

门控消融实验的结果（Gramian门控在所有测试场景下与无门控的RMSE差异$<2\%$）是本文最意外的发现。这促使我们重新审视PC-RCO的鲁棒性来源。分析表明，以下三个结构性特征共同保证了失配下的稳定性：

\begin{enumerate}
    \item \textbf{速度层预测-校正}：速度预测误差天然仅为$O(10^{-4})$ m/s量级（稳态时推力与阻力近似平衡，海流对阻力影响为二阶小量），限制了更新信号的幅度。
    \item \textbf{GM时间衰减}：$\tau_c=100$s的持续衰减将估计拉向先验均值，阻止了长时间漂移。
    \item \textbf{单步限幅}：$\Delta c_{\max}=0.1$ m/s确保任何单次异常更新不会造成灾难性后果。
\end{enumerate}

对比之下，CFD-Luenberger观测器（朴素模型基实现，缺乏限幅和GM衰减）在直线轨迹崩塌至RMSE$_{V_c}=2.284$，为上述分析提供了反面证据。需要指出的是，本文未做逐机制消融，因此上述分析应理解为"三机制作为整体的鲁棒性"，而非每个机制的独立因果贡献。

\subsection{与概率方法的对比}

EKF在无失配时精度最高(RMSE$_{V_c}=0.035$)，但对模型失配极度敏感(退化77\%)。PC-RCO在精度(0.064)和鲁棒性(退化42\%)之间取得了工程上有意义的折中。从部署角度看，PC-RCO的优势在于：
\begin{itemize}
    \item 状态维度低(2 vs EKF的8)，计算更轻
    \item 无需EKF式的协方差矩阵设计与调优
    \item 各参数具有直观的物理含义（$\tau_c$为海流记忆时间，$\Delta c_{\max}$为最大变化率）
\end{itemize}

\subsection{激励门控的反思}

本文最初的设计意图是用Gramian门控防止弱激励下的错误更新。然而消融实验表明，在DVL测量噪声$\sigma_v=0.02$ m/s条件下，$\lambda_{\min}(\bm{W}_c)$始终高于门控阈值$10^{-8}$，硬门控从未实际触发。这一结果并不意味着Gramian分析无价值，而是提示：在工程实践中，连续的正则化（通过Gramian条件数调节有效增益）比二值门控更实用。Gramian分析的价值因此重新定位为\textbf{可观测性诊断与自适应正则化}工具。

\subsection{实际部署考量}

PC-RCO的计算复杂度可控：每次更新需要3次模型前向计算（基准+2次数值扰动），在XHY AUV的嵌入式平台上的典型计算时间$<3$ ms，满足实时性要求。DVL短暂失效时，GM传播可维持海流估计约30--60秒。若DVL长时间失效，估计将退化至先验均值，需重新获取DVL锁定后的机动激励以恢复精度。
